\chapter{Conception Globale du Syst\`eme}
\minitoc


%--------------------------------------------------------------
\section{Introduction}
%--------------------------------------------------------------

\todo{Objectif du chapitre : vision complète du système avant les détails
d'implémentation.}

%--------------------------------------------------------------
\section{Architecture générale du système}
\label{sec:archi_globale}
%--------------------------------------------------------------

\todo{%
  Figure PLEINE PAGE obligatoire : tous les services (boîtes),
  tous les protocoles (WebSocket, REST, ESL, SIP/RTP),
  les ports, une légende, une frontière "données internes / données externes".
}

\begin{figure}[H]
  \centering
  \input{tikz/fig3_architecture_globale}
  \caption{Architecture globale du système}
  \label{fig:architecture_globale}
\end{figure}

% Tableau de référence des services
\begin{longtable}{lllll}
  \caption{Tableau des services du système} \label{tab:services} \\
  \toprule
  \textbf{Service} & \textbf{Technologie} & \textbf{Port} & \textbf{Rôle} & \textbf{Communication} \\
  \midrule
  \endfirsthead
  \multicolumn{5}{c}{\textit{(Suite du Tableau~\ref{tab:services})}} \\
  \toprule
  \textbf{Service} & \textbf{Technologie} & \textbf{Port} & \textbf{Rôle} & \textbf{Communication} \\
  \midrule
  \endhead
  \bottomrule
  \endfoot
  FreeSWITCH   & FreeSWITCH 1.10.11 & SIP:5060 / ESL:8022 & Serveur SIP/RTP   & WS $\rightarrow$ Voice GW \\
  Voice Gateway & FastAPI            & 8000                & STT, TTS, VAD     & WS $\leftarrow$ FS, REST $\rightarrow$ Rasa \\
  Rasa NLU/Core & Rasa 3.6.x        & 5005                & NLU, dialogue     & REST webhook \\
  Action Server & FastAPI custom     & 5055                & Actions, LLM, RAG & Appelé par Rasa \\
  KB Service   & FastAPI            & 8001                & ChromaDB, RAG     & REST \\
  Ollama       & Ollama             & 11434               & Embeddings bge-m3 & REST \\
  EspoCRM      & PHP/Apache         & 80                  & CRM débiteurs     & REST \\
  Duckling     & Duckling           & 8080                & Entités temporelles & REST \\
  Django Backend & Django + Gunicorn & 8080 (ext.)        & API gouvernance   & REST $\leftarrow$ Next.js \\
  Celery Worker & Celery 5.4        & —                   & Tâches async      & Redis broker \\
  Redis        & Redis              & 6379                & Broker, cache     & — \\
  PostgreSQL   & PostgreSQL         & 5432                & DB Django         & — \\
  Next.js Frontend & Next.js        & 3000                & Portail           & REST $\rightarrow$ Django \\
\end{longtable}

%--------------------------------------------------------------
\section{Flux de données principal : un appel de recouvrement de bout en bout}
%--------------------------------------------------------------

\subsection{Narration du scénario nominal}

\todo{%
  Raconter le parcours d'un appel étape par étape :
  1. Déclenchement Celery/ESL
  2. FreeSWITCH compose le numéro
  3. mod\_audio\_stream $\rightarrow$ Voice Gateway
  4. VAD détecte l'activité vocale
  5. STT transcrit l'utterance
  6. POST texte $\rightarrow$ Rasa
  7. NLU classifie l'intent
  8. Action Server : EspoCRM / ChromaDB / Cerebras / DistilBERT (async)
  9. Réponse $\rightarrow$ Piper TTS $\rightarrow$ resampling $\rightarrow$ ESL uuid\_broadcast
}

\subsection{Diagramme de séquence — scénario nominal}

% ---- DIAGRAMME DE SÉQUENCE PRINCIPAL (pgf-umlsd) ----
\begin{figure}[H]
  \centering
  \begin{sequencediagram}
    \newthread{caller}{Débiteur}
    \newinst[1]{fs}{FreeSWITCH}
    \newinst[1]{vg}{Voice Gateway}
    \newinst[1]{rasa}{Rasa}
    \newinst[1]{action}{Action Server}
    \newinst[1]{crm}{EspoCRM}
    \newinst[1]{tts}{Piper TTS}

    \begin{call}{caller}{Appel SIP entrant}{fs}{Connexion établie}
    \end{call}

    \begin{call}{fs}{Audio PCM 8kHz (WebSocket)}{vg}{}
      \begin{call}{vg}{VAD $\rightarrow$ STT}{vg}{Texte transcrit}
      \end{call}
      \begin{call}{vg}{POST /webhooks/rest/webhook}{rasa}{Réponse texte}
        \begin{call}{rasa}{Exécuter action}{action}{Résultat}
          \begin{call}{action}{GET /api/v1/Contact}{crm}{Données débiteur}
          \end{call}
        \end{call}
      \end{call}
      \begin{call}{vg}{Synthèse audio}{tts}{PCM 22050Hz}
      \end{call}
    \end{call}

    \begin{call}{fs}{ESL uuid\_broadcast}{caller}{Audio joué}
    \end{call}

  \end{sequencediagram}
  \caption{Diagramme de séquence — Scénario nominal d'un appel de recouvrement}
  \label{fig:seq_nominal}
\end{figure}

\subsection{Diagramme de séquence — escalade vers agent humain}

\begin{figure}[H]
  \centering
  \input{tikz/fig3_seq_escalade}
  \caption{Diagramme de séquence — Escalade vers agent humain}
  \label{fig:seq_escalade}
\end{figure}

%--------------------------------------------------------------
\section{Acteurs et cas d'utilisation}
%--------------------------------------------------------------

\subsection{Acteurs identifiés}

\todo{Débiteur, Superviseur bancaire, Administrateur système, Agent vocal IA.}

\subsection{Diagramme de cas d'utilisation global}

\begin{figure}[H]
  \centering
  \input{tikz/fig3_usecase}
  \caption{Diagramme de cas d'utilisation global}
  \label{fig:usecase}
\end{figure}

\subsection{Description des cas d'utilisation principaux}

\begin{longtable}{lllp{5cm}}
  \caption{Description des cas d'utilisation principaux} \label{tab:usecase} \\
  \toprule
  \textbf{UC} & \textbf{Acteur} & \textbf{Description} & \textbf{Préconditions} \\
  \midrule
  \endfirsthead
  \toprule
  \textbf{UC} & \textbf{Acteur} & \textbf{Description} & \textbf{Préconditions} \\
  \midrule
  \endhead
  \bottomrule
  \endfoot
  UC01 & Débiteur       & Recevoir un appel de recouvrement       & Dossier actif dans EspoCRM \\
  UC02 & Débiteur       & S'authentifier (code postal / DDN)      & Appel en cours \\
  UC03 & Débiteur       & Négocier un échéancier de paiement      & Authentifié \\
  UC04 & Débiteur       & Être transféré à un agent humain        & Détresse émotionnelle \\
  UC05 & Superviseur    & Consulter le monitoring des appels      & Connecté au portail \\
  UC06 & Superviseur    & Consulter les analytics                 & Connecté au portail \\
  UC07 & Administrateur & Gérer les guardrails IA                 & Rôle admin \\
  UC08 & Administrateur & Lancer une campagne d'appels            & Rôle admin \\
  UC09 & Administrateur & Entraîner le modèle Rasa (Forge)        & Rôle admin \\
  UC10 & Administrateur & Exporter les logs réglementaires        & Rôle admin \\
\end{longtable}

%--------------------------------------------------------------
\section{Modèle de domaine}
%--------------------------------------------------------------

\begin{figure}[H]
  \centering
  \input{tikz/fig3_class_diagram}
  \caption{Diagramme de classes — Modèle de domaine}
  \label{fig:class_diagram}
\end{figure}

\todo{%
  Classes : Contact/Debtor, Call, Utterance, EmotionLog,
  Experiment, Run, Guardrail, KBDocument, Campaign, AuditLog.
}

%--------------------------------------------------------------
\section{Diagramme de déploiement}
%--------------------------------------------------------------

\todo{%
  Distinct de l'architecture globale (Section~\ref{sec:archi_globale}).
  Montrer les nœuds physiques/virtuels : machine WSL2 (hôte Windows),
  réseau Docker interne, réseau SIP, accès externe Cerebras.
  Pour chaque nœud : artefacts déployés, gestionnaire de processus, protocole réseau.
}

\begin{figure}[H]
  \centering
  \input{tikz/fig3_deployment}
  \caption{Diagramme de déploiement — Infrastructure physique et virtuelle}
  \label{fig:deployment}
\end{figure}

%--------------------------------------------------------------
\section{Maquettes IHM du portail de gouvernance}
%--------------------------------------------------------------

\todo{%
  Présenter les maquettes (wireframes) des pages principales du portail
  \textbf{avant} la réalisation (Ch.4). Ces maquettes définissent la cible
  fonctionnelle et visuelle validée avec l'encadrant entreprise.
  Outil recommandé : Figma, Balsamiq, ou croquis annoté scanné.
  Un wireframe n'a pas besoin d'être pixel-perfect — l'essentiel est de montrer
  la disposition des éléments et les interactions principales.
}

\subsection{Vue d'ensemble — navigation du portail}

\begin{figure}[H]
  \centering
  \input{tikz/fig3_ihm_navigation}
  \caption{Maquette IHM — Structure de navigation du portail}
  \label{fig:ihm_navigation}
\end{figure}

\subsection{Page Agent Dashboard}

\begin{figure}[H]
  \centering
  % \includegraphics[width=\textwidth]{figures/wireframe_dashboard.png}
  \todo{Maquette : Agent Dashboard — KPIs, appels actifs, alertes émotion}
  \caption{Maquette IHM — Agent Dashboard}
  \label{fig:wireframe_dashboard}
\end{figure}

\subsection{Page Forge — Experiments}

\begin{figure}[H]
  \centering
  % \includegraphics[width=\textwidth]{figures/wireframe_forge.png}
  \todo{Maquette : Forge Experiments — tableau runs, métriques NLU, bouton Train}
  \caption{Maquette IHM — Forge Experiments}
  \label{fig:wireframe_forge}
\end{figure}

\subsection{Page Guardrails}

\begin{figure}[H]
  \centering
  % \includegraphics[width=\textwidth]{figures/wireframe_guardrails.png}
  \todo{Maquette : Guardrails — liste règles actives, modal création}
  \caption{Maquette IHM — Guardrails}
  \label{fig:wireframe_guardrails}
\end{figure}

%--------------------------------------------------------------
\section{Sécurité et conformité}
%--------------------------------------------------------------

\subsection{Flux d'authentification 2FA $\rightarrow$ JWT}

\begin{figure}[H]
  \centering
  \begin{sequencediagram}
    \newthread{user}{Utilisateur}
    \newinst[1]{nextjs}{Next.js}
    \newinst[1]{django}{Django}
    \newinst[1]{email}{Serveur SMTP}

    \begin{call}{user}{email + mot de passe}{nextjs}{}
      \begin{call}{nextjs}{POST /auth/login/}{django}{OTP généré}
        \begin{call}{django}{Envoi OTP}{email}{}
        \end{call}
      \end{call}
    \end{call}

    \begin{call}{user}{Saisie OTP}{nextjs}{}
      \begin{call}{nextjs}{POST /auth/verify-otp/}{django}{access\_token + refresh\_token}
      \end{call}
    \end{call}

    \begin{call}{user}{Requête API}{nextjs}{}
      \begin{call}{nextjs}{Bearer token}{django}{Réponse}
      \end{call}
    \end{call}

  \end{sequencediagram}
  \caption{Diagramme de séquence — Authentification 2FA $\rightarrow$ JWT}
  \label{fig:seq_auth}
\end{figure}

\subsection{Cartographie de résidence des données}

\begin{table}[H]
  \centering
  \caption{Cartographie des flux de données et résidence}
  \label{tab:residence_donnees}
  \begin{tabularx}{\textwidth}{lXll}
    \toprule
    \textbf{Flux} & \textbf{Données transmises} & \textbf{Périmètre} & \textbf{Justification} \\
    \midrule
    Voice GW $\rightarrow$ Rasa    & Texte transcrit          & Interne  & Réseau local \\
    Action Server $\rightarrow$ EspoCRM & Données débiteur    & Interne  & Réseau local \\
    Action Server $\rightarrow$ ChromaDB & Query texte        & Interne  & Process local \\
    \rowcolor{yellow!20}
    Action Server $\rightarrow$ \textbf{Cerebras} & Texte + contexte \ac{RAG} & \textbf{Externe} & Texte générique, \textbf{aucun PII} \\
    Tous autres flux & Texte ou audio              & Interne  & Réseau local \\
    \bottomrule
  \end{tabularx}
\end{table}

%--------------------------------------------------------------
\section{Conclusion}
%--------------------------------------------------------------

\todo{%
  Synthèse de la conception. Transition : \og{}La conception globale
  étant établie, le chapitre suivant détaille l'implémentation de chacun
  des 7 modules du système.\fg{}
}
